% Two-way realization taxonomy: replaces the Preserved / Carrier-transferred /
% Rematerialized itemize in \S\ref{sec3.3} (root.tex "Realization strategy").

\paragraph{Realization strategy}
Transferring a gesture to a different morphology changes which carriers and attributes are physically available, and a target body affords two qualitatively different responses. Movement structure is comparatively portable: an axis role, a signed stroke, or a repetition pattern can be reproduced by whatever carrier the target body provides, even when that carrier is anatomically unrelated to the original one. Orientation and body-relative location are not portable in the same way, because they depend on landmarks and reachable regions that may simply not exist. We therefore distinguish two strategies by \emph{what} is carried over, rather than by which body part carries it.

\begin{itemize}
\item \textbf{Retargeted:} The recognition-critical movement of the source gesture is applied, essentially unchanged, to a carrier the target body already has. The stroke keeps its signed structure and repetition, and is re-expressed in the new carrier's own frame; only the carrier changes. A head nod becomes a nod of TIAGo's head, an oscillation of a single arm's wrist, or a fore--aft rock of a mobile base, and \texttt{gaze\_shift\_avoid} becomes a turn of whichever segment establishes facing. Cues whose identity rests mainly on a movement axis, and cues that a single arm can execute directly, are typically realizable by retargeting alone.

\item \textbf{Rematerialized:} No available carrier can express the source movement as it stands, so the communicative structure itself is rebuilt. The realization modifies the movement, expresses the meaning through a different structure, or recruits additional evidence that the source did not use: bilateral arm flapping becomes a bird-like trajectory traced by a planar base; a beckoning wrist flap becomes an approach-and-retreat of the whole body; a raised arm is accompanied by head scanning so that the intended reading survives. Rematerialization is thus a change of communicative form, not a change of body part.
\end{itemize}

The distinction matters because the two strategies fail differently. A retargeted gesture inherits its source evidence and fails mainly when an attribute is instantiated imprecisely. A rematerialized gesture has no source form to fall back on: it must reconstruct enough evidence to remain distinguishable from its neighbours, and its characteristic failure is that several meanings converge on the same remaining motion.
